\documentclass[a4paper,12pt]{article}

\usepackage[utf8]{inputenc}
\usepackage[T1]{fontenc}
\usepackage[brazil]{babel}
\usepackage{amsmath}

\title{Funcionamento da Seed e Geração de Números Pseudoaleatórios}
\author{}
\date{}

\begin{document}

\maketitle

\section{O que é uma Seed?}

Uma \textit{seed} (semente) é o valor inicial utilizado por um gerador de números pseudoaleatórios.
Quando a mesma seed é utilizada, a sequência de números gerada será sempre a mesma.

Por exemplo:

\begin{verbatim}
Random random = new Random(42);
\end{verbatim}

Toda execução iniciada com a seed 42 produzirá a mesma sequência de números.

\section{Gerador Linear Congruente}

O algoritmo clássico utilizado em diversos geradores pseudoaleatórios é o
\textit{Linear Congruential Generator} (LCG), definido pela equação:

\[
X_{n+1} = (aX_n + c) \bmod m
\]

onde:

\begin{itemize}
    \item $X_n$ é o estado atual;
    \item $a$ é o multiplicador;
    \item $c$ é o incremento;
    \item $m$ é o módulo;
    \item $X_{n+1}$ é o próximo valor gerado.
\end{itemize}

\section{Exemplo Simples}

Considere o gerador fictício:

\[
X_{n+1} = (3X_n + 7) \bmod 100
\]

com seed inicial:

\[
X_0 = 42
\]

\subsection*{Primeiro número}

\[
X_1 = (3 \times 42 + 7) \bmod 100
\]

\[
X_1 = 133 \bmod 100
\]

\[
X_1 = 33
\]

\subsection*{Segundo número}

\[
X_2 = (3 \times 33 + 7) \bmod 100
\]

\[
X_2 = 106 \bmod 100
\]

\[
X_2 = 6
\]

\subsection*{Terceiro número}

\[
X_3 = (3 \times 6 + 7) \bmod 100
\]

\[
X_3 = 25
\]

A sequência produzida é:

\[
42 \rightarrow 33 \rightarrow 6 \rightarrow 25 \rightarrow \cdots
\]

Sempre que a seed inicial for 42, a mesma sequência será obtida.

\section{Aplicação no Experimento}

Nos testes dos algoritmos de ordenação, a utilização de uma seed fixa garante
que todos os algoritmos recebam exatamente o mesmo conjunto de dados,
permitindo uma comparação justa dos tempos de execução.

\end{document}